Table~\ref{tab:main} gives the main results (mean $\pm$ std over \rhval{const/n_seeds} seeds, \rhval{const/n_probs} problems per task and seed) and Fig.~\ref{fig:perbin} the per-bin solved rates, including two ablated variants discussed in Sec.~\ref{sec:ablations}. Paired statistics come from \texttt{rh compare} (paired $t$-test over seeds).

\begin{table}[t]\centering\caption{Main results: solved rate within \rhval{const/budget} expansions and mean expansions (failures count as the budget). \emph{val} is the training size range, the \emph{test} tasks are larger (connective ranges in the task name). Rows for the tuning bin \emph{tune\_21\_35} omit the heuristic (its rows are in Table~\ref{tab:tune}) and are shown only for completeness; bold/underline in this table ignore the missing entry and are not meaningful there.}\label{tab:main}
\resizebox{\columnwidth}{!}{\begin{tabular}{llccc}
\toprule
Method & Task & solved $\uparrow$ & mean\_exp $\downarrow$ & $n$ \\
\midrule
BFS & val\_3\_10 & 0.91 $\pm$ 0.03 & 21.03 $\pm$ 3.42 & 5 \\
Hand heuristic & val\_3\_10 & \textbf{1.00 $\pm$ 0.00} & \textbf{4.42 $\pm$ 0.28} & 5 \\
Feature-MLP policy & val\_3\_10 & \underline{1.00 $\pm$ 0.00} & \underline{5.26 $\pm$ 0.76} & 5 \\
\textbf{Transformer policy} & val\_3\_10 & 1.00 $\pm$ 0.00 & 5.99 $\pm$ 0.83 & 5 \\
\midrule
BFS & test\_11\_20 & 0.47 $\pm$ 0.03 & 63.33 $\pm$ 2.73 & 5 \\
Hand heuristic & test\_11\_20 & \textbf{1.00 $\pm$ 0.00} & \textbf{11.31 $\pm$ 0.45} & 5 \\
Feature-MLP policy & test\_11\_20 & \underline{0.89 $\pm$ 0.02} & \underline{27.16 $\pm$ 1.86} & 5 \\
\textbf{Transformer policy} & test\_11\_20 & 0.82 $\pm$ 0.02 & 34.99 $\pm$ 1.37 & 5 \\
\midrule
BFS & test\_21\_40 & 0.27 $\pm$ 0.06 & 80.82 $\pm$ 5.00 & 5 \\
Hand heuristic & test\_21\_40 & \textbf{0.96 $\pm$ 0.01} & \textbf{27.57 $\pm$ 2.35} & 5 \\
Feature-MLP policy & test\_21\_40 & \underline{0.58 $\pm$ 0.03} & \underline{55.59 $\pm$ 3.83} & 5 \\
\textbf{Transformer policy} & test\_21\_40 & 0.52 $\pm$ 0.05 & 59.82 $\pm$ 5.14 & 5 \\
\midrule
BFS & test\_41\_70 & 0.20 $\pm$ 0.04 & 85.69 $\pm$ 1.69 & 5 \\
Hand heuristic & test\_41\_70 & \textbf{0.73 $\pm$ 0.03} & \textbf{47.34 $\pm$ 2.38} & 5 \\
Feature-MLP policy & test\_41\_70 & \underline{0.46 $\pm$ 0.03} & \underline{65.25 $\pm$ 3.20} & 5 \\
\textbf{Transformer policy} & test\_41\_70 & 0.41 $\pm$ 0.05 & 68.81 $\pm$ 4.19 & 5 \\
\midrule
BFS & tune\_21\_35 & 0.30 $\pm$ 0.04 & 78.79 $\pm$ 3.25 & 5 \\
Feature-MLP policy & tune\_21\_35 & \textbf{0.63 $\pm$ 0.06} & \textbf{51.56 $\pm$ 4.28} & 5 \\
\textbf{Transformer policy} & tune\_21\_35 & \underline{0.55 $\pm$ 0.05} & \underline{57.31 $\pm$ 3.81} & 5 \\
\bottomrule
\end{tabular}
}
\end{table}
\begin{figure}[t]\centering\includegraphics[width=0.95\columnwidth]{perbin_fig.pdf}\caption{Solved rate per test-size bin (mean, std over seeds). ``Transf.'' and ``MLP'' use $-\log p$ cost; ``Transf. greedy'' follows the transformer without backtracking; ``MLP rank'' uses the MLP's rank as the step cost (post-hoc variants).}\label{fig:perbin}
\end{figure}

\paragraph{In the training size range} all guided systems are essentially perfect (solved: transformer \rhval{main/transformer-policy/val_3_10/solved/mean:3}, MLP \rhval{main/feature-mlp-policy/val_3_10/solved/mean:3}, heuristic \rhval{main/hand-heuristic/val_3_10/solved/mean:3}) while BFS already misses part of the problems (\rhval{main/bfs/val_3_10/solved/mean:2}); on this bin the heuristic's advantage over the transformer is not significant (paired $p=$ \rhval{cmp/main/hand-heuristic/val_3_10/solved/paired_p}). The held-out top-1 accuracy against the optimal tactic set is \rhval{main/transformer-policy/val_3_10/val_top1/mean:3} for the transformer and \rhval{main/feature-mlp-policy/val_3_10/val_top1/mean:3} for the MLP, i.e.\ the transformer imitates optimal tactics less accurately than the MLP given the same number of steps.

\paragraph{H1 (policy $>$ BFS): supported.} In all three test bins the transformer solves more problems than BFS: at 11--20 (\rhval{main/transformer-policy/test_11_20/solved/mean:2} vs.\ \rhval{main/bfs/test_11_20/solved/mean:2}, $p=$ \rhval{cmp/main/bfs/test_11_20/solved/paired_p}), at 21--40 (\rhval{main/transformer-policy/test_21_40/solved/mean:2} vs.\ \rhval{main/bfs/test_21_40/solved/mean:2}, $p=$ \rhval{cmp/main/bfs/test_21_40/solved/paired_p}) and at 41--70 (\rhval{main/transformer-policy/test_41_70/solved/mean:2} vs.\ \rhval{main/bfs/test_41_70/solved/mean:2}, $p=$ \rhval{cmp/main/bfs/test_41_70/solved/paired_p}). BFS is a weak reference: it ignores the structure of the problem and its expansion count is dominated by the branching factor.

\paragraph{H2 (policy $>$ heuristic at 41--70): refuted.} The tuned heuristic solves \rhval{main/hand-heuristic/test_41_70/solved/mean:2} against \rhval{main/transformer-policy/test_41_70/solved/mean:2} for the transformer (difference \rhval{cmp/main/hand-heuristic/test_41_70/solved/delta}, paired $p=$ \rhval{cmp/main/hand-heuristic/test_41_70/solved/paired_p}; the \texttt{delta} is transformer minus heuristic, the sign convention of \texttt{rh compare}) and \rhval{main/feature-mlp-policy/test_41_70/solved/mean:2} for the feature MLP. The heuristic also needs fewer expansions on average (\rhval{main/hand-heuristic/test_41_70/mean_exp/mean:1} vs.\ \rhval{main/transformer-policy/test_41_70/mean_exp/mean:1}). The heuristic is not perfect: it still fails on a large fraction of the 41--70 problems, i.e.\ these are cases where the budget, not the ordering of non-branching rules, binds.

\paragraph{H3 (gap shrinks with size): refuted.} The transformer-minus-heuristic gap in solved rate is \rhval{analysis/analysis/test_41_70/gap_tf_heur_t11/mean:2} (11--20), \rhval{analysis/analysis/test_41_70/gap_tf_heur_t21/mean:2} (21--40) and \rhval{analysis/analysis/test_41_70/gap_tf_heur_t41/mean:2} (41--70). The gap is not monotone: it is smallest at 11--20, largest at 21--40, and narrower at 41--70. Part of the narrowing at 41--70 is a ceiling effect, because the heuristic itself drops to \rhval{main/hand-heuristic/test_41_70/solved/mean:2}. We therefore cannot claim that the policy degrades \emph{faster} than the heuristic with size; we can only say the policy is behind in every bin. (H3 was formulated as a monotone decrease of the policy-minus-heuristic difference, so the pre-registered criterion is not met.)

\paragraph{Where the transformer fails.} The transformer is below the MLP in every test bin (paired differences MLP minus transformer: \rhval{analysis/analysis/test_41_70/d_mlp_tf_t11/mean:2}, \rhval{analysis/analysis/test_41_70/d_mlp_tf_t21/mean:2}, \rhval{analysis/analysis/test_41_70/d_mlp_tf_t41/mean:2}; $p=$ \rhval{analysis/analysis/test_41_70/p_mlp_tf_t11/mean}, \rhval{analysis/analysis/test_41_70/p_mlp_tf_t21/mean}, \rhval{analysis/analysis/test_41_70/p_mlp_tf_t41/mean}), although the MLP sees only hand-chosen features. The sequence model has lower validation accuracy and, in the largest bin, also the largest variance between seeds (std \rhval{main/transformer-policy/test_41_70/solved/std:2}). Per-regime numbers in Table~\ref{tab:main} show that the failure is a size effect for all learned systems: their solved rate falls from the training range to 41--70 much faster than the heuristic's does at 11--20 and 21--40.

\paragraph{Diagnostic: do the policies defer branching rules?} The heuristic's central principle is to apply non-branching tactics before branching ones. On the states visited by the greedy heuristic proof of each test problem we measured how often the policy's top-1 tactic is non-branching among states that offer both kinds (\emph{defer rate}) and how often it equals the heuristic's top-1 tactic (Table~\ref{tab:diag}). The transformer's defer rate stays high in all bins, while the MLP's falls strongly with size; the two policies' top-1 agreement with the heuristic is only moderate even in the training range (which is expected, because many tactics tie in optimal size). These measurements describe the policies' behaviour; they do not by themselves show that deferring branching explains the search outcomes.

\begin{table}[t]\centering\caption{Diagnostic on heuristic-proof states: top-1 agreement with the heuristic and defer rate (non-branching preferred when both kinds are available). Mean $\pm$ std over seeds; bold/underline marks in this table are not meaningful.}\label{tab:diag}
\resizebox{\columnwidth}{!}{\begin{tabular}{llccc}
\toprule
Method & Task & agree\_heur $\downarrow$ & defer\_rate $\downarrow$ & $n$ \\
\midrule
\textbf{Transformer policy} & val\_3\_10 & \underline{0.72 $\pm$ 0.02} & \underline{0.93 $\pm$ 0.01} & 5 \\
Feature-MLP policy & val\_3\_10 & \textbf{0.69 $\pm$ 0.03} & \textbf{0.83 $\pm$ 0.02} & 5 \\
\midrule
\textbf{Transformer policy} & test\_11\_20 & \underline{0.64 $\pm$ 0.02} & \underline{0.90 $\pm$ 0.02} & 5 \\
Feature-MLP policy & test\_11\_20 & \textbf{0.58 $\pm$ 0.00} & \textbf{0.80 $\pm$ 0.01} & 5 \\
\midrule
\textbf{Transformer policy} & test\_21\_40 & \underline{0.58 $\pm$ 0.02} & \underline{0.86 $\pm$ 0.04} & 5 \\
Feature-MLP policy & test\_21\_40 & \textbf{0.48 $\pm$ 0.01} & \textbf{0.72 $\pm$ 0.02} & 5 \\
\midrule
\textbf{Transformer policy} & test\_41\_70 & \underline{0.55 $\pm$ 0.02} & \underline{0.84 $\pm$ 0.04} & 5 \\
Feature-MLP policy & test\_41\_70 & \textbf{0.40 $\pm$ 0.03} & \textbf{0.61 $\pm$ 0.04} & 5 \\
\bottomrule
\end{tabular}
}
\end{table}
